\section{Sizes, a lower bound and a depth-two example for stars}\label{app:star}

For a rational $r=p/q$ in lowest terms let $\operatorname{ht}(r)=\log_2\max\{\abs p,q\}$.
Then $\operatorname{ht}(rs),\operatorname{ht}(r/s)\le\operatorname{ht}(r)+\operatorname{ht}(s)$
and $\operatorname{ht}(r_1+\dots+r_N)\le\log_2N+\sum_j\operatorname{ht}(r_j)$. Let
$\tau_i$ bound $\operatorname{ht}$ over the data of leaf $i$ (its coefficients,
bounds and rows), $\tau_0$ over the center data (its coefficients, bounds and
the rows in $y$ alone), and $\tau=\max_{0\le i\le k}\tau_i$.

\paragraph{Sizes.}
The coefficients of the affine functions defining $L_i$ and $U_i$ have
$\operatorname{ht}\le2\tau_i$, and $v_i$ has $\operatorname{ht}\le2\tau_i+1$. A
breakpoint is the intersection of two such functions or the zero of the affine
factor in~\eqref{eq:endpoint-factor}; in both cases $\operatorname{ht}\le12\tau+4$,
independently of $k$. The endpoints of $I$ are bounds of $y$, zeros of rows in
$y$ alone ($\operatorname{ht}\le2\tau_0$) or intersections of two lines of one
leaf, so they obey the same bound. On a piece, the contribution of leaf $i$ to
$V$ is a quadratic in $y$ whose coefficients have $\operatorname{ht}\le11\tau_i+7$,
so the coefficients of $V$, which are sums of $k+1$ such terms, have
$\operatorname{ht}\le C:=\log_2(k+1)+\tau_0+\sum_i(11\tau_i+7)$, which is linear in
the input length. Evaluating at a breakpoint, at an endpoint of $I$ or at a
stationary point $-c_1/(2c_2)$ gives a minimum with $\operatorname{ht}\le4C+36\tau+14$.
Because the sweep updates the coefficients exactly, every running sum equals
the coefficient of the current piece and obeys the same bound. The algorithm
is therefore strongly polynomial.

\paragraph{Lower bound.}
Ben-Or's theorem states that an algebraic computation tree deciding membership
in a set $W\subseteq\R^n$ with $N$ connected components has depth
$\Omega(\log N-n)$ \citep{BenOr1983}. For $r\ge1$ let
$W=\{z\in(0,2r)^r:\abs{z_i-z_j}\ge1\ \text{for}\ i\ne j\}$, which has one component
for each of the $r!$ orderings. For $z\in(0,2r)^r$ build a star with center
$y\in[0,2r]$ and, for each $j$, three leaves: $x\in[0,1]$ with objective
$(y-z_j+1)x$; $x\in[0,1]$ with objective $(z_j+1-y)x$; and $x\in[0,2r+1]$ with
the row $x\ge y-z_j$ and objective $2x$; add the center terms
$-ry+\sum_jz_j-r$. With $t=y-z_j$, the identity
$\min\{0,t+1\}+\min\{0,1-t\}+2\max\{0,t\}-t-1=-\max\{0,1-\abs t\}$ shows that
$V(y)=-\sum_j\max\{0,1-\abs{y-z_j}\}$, a sum of negative tents of width two. If
$z\in W$, at most two tents are positive at any $y$ and their sum is at most
one, so $\min V\ge-1$; if $\abs{z_i-z_j}<1$, then $V<-1$ at the midpoint of $z_i$
and $z_j$. The instance has $3r$ leaves and $r$ leaf rows and is computed from
$z$ with $O(r)$ operations, so deciding $\min V\ge-1$ needs depth
$\Omega(r\log r)$.

\paragraph{Depth two.}
For $f=(x_2-\tfrac18)x_1+x_2^2-x_2+x_2x_3+x_3^2-x_3$ on $[0,1]^3$, minimizing over
$x_1$ first gives $\min\{0,x_2-\tfrac18\}+x_2^2+(t-1)x_2$ with $t=x_3$. On
$[0,\tfrac18]$ its minimum is $-\tfrac18$ at $x_2=0$; on $[\tfrac18,1]$ it is
$-\tfrac14(1-t)^2$ at $x_2=(1-t)/2$ when $t\le\tfrac34$, and larger than $-\tfrac18$
otherwise. Hence the conditional value is $\min\{-\tfrac18,-\tfrac14(1-t)^2\}$,
whose two pieces meet at $t=1-\sqrt2/2$, a root of $2t^2-4t+1$. Adding $t^2-t$
and minimizing over $t\in[0,1]$ gives $\min f=-\tfrac38$ at $(1,0,\tfrac12)$.
